\section{Marco Teórico}
\label{ch:marcoteroico}

\subsection{Modernización de Sistemas}
\label{sec:modernizacion_sistemas}
Actualmente, la tecnología avanza tan rápido que las organizaciones e instituciones no se pueden quedar atrás. Por lo tanto, actualizar los sistemas de software se ha vuelto una de las tareas más importantes para los directores y encargados de la tecnología en las empresas. El "modernizar un sistema", no se refiere a hacer un cambio de la noche a la mañana, sino a un proceso continuo de mejora para que el software no pierda su valor. Según explican los investigadores Abu Bakar, Razali y Jambari (2020), modernizar consiste en renovar esas aplicaciones viejas utilizando tecnologías actuales o estructuras más modernas. El objetivo de esto es que el sistema pueda reaccionar rápido a los cambios que necesita la organización, funcione mejor y sea mucho más seguro.\cite{AbuBakar2020LegacyModernization}
\subsubsection{Sistemas Heredados (Legacy Systems)}
\label{sec:sistemas_heredados}
Un concepto clave en esta área son los llamados "sistemas heredados" o legacy systems. Básicamente, son sistemas de información viejos que todavía se usan porque son el motor principal de las operaciones diarias de una empresa, pero que ya se quedaron atrasados respecto a las tecnologías y estándares actuales. De acuerdo con Abu Bakar, Razali y Jambari (2020), estos sistemas utilizan lenguajes de programación y herramientas de generaciones pasadas.\cite{AbuBakar2020LegacyModernization} Aunque suelen dar problemas técnicos o fallas constantes, las organizaciones no los pueden dejar de usar tan fácilmente porque tienen un valor muy alto. Esto pasa porque el sistema viejo guarda las reglas del negocio, los datos históricos de muchos años y la costumbre de los usuarios que ya se la saben de memoria; apagar el sistema de golpe afectaría por completo el trabajo de todos los días.

El verdadero problema es que mantener un sistema viejo con el tiempo se vuelve una pesadilla muy cara y lenta. Cada vez es más difícil encontrar programadores que conozcan herramientas antiguas, y por lo general, la documentación o los manuales para guiarse ya no existen o están incompletos. Al final, el sistema se vuelve una barrera que no deja que la institución innove o agregue nuevas funciones que los usuarios necesitan.

Sobre esto, Mandviwalla, Tiwana y Bradshaw (2026) comentan que los sistemas viejos y personalizados terminan atrapando los datos en diferentes bases de datos que no se hablan entre sí. Intentar arreglar esto con soluciones rápidas o provisionales solo aumenta la "deuda técnica" (el desorden en el código), creando una trampa que eleva el riesgo de que el sistema se caiga o sufra ataques cibernéticos.\cite{mandviwalla2025modernizing} Por lo tanto, la mejor opción es aplicar técnicas de ingeniería inversa (revisar el código viejo para entender qué hace) y migrar la plataforma hacia un entorno moderno, recuperando la flexibilidad del software pero cuidando la información de la institución.

\subsection{Evolución de las Tecnologías de Desarrollo Web en Microsoft}
\label{sec:evolucion_tecnologias}
El desarrollo de aplicaciones web con Microsoft ha pasado por una gran transformación a lo largo de las últimas décadas. Lo que comenzó como un modelo pensado para servidores locales y sistemas operativos Windows, ha evolucionado hacia un entorno completamente abierto y diseñado para la nube. Esta evolución no solo representa un cambio de herramientas, sino un cambio completo en la manera de programar de los desarrolladores y en las metodologías usadas por las organizaciones para actualizar sus aplicaciones.

\subsubsection{La Era de ASP.NET Web Forms}
\label{sec:aspnet_webforms}
Lanzado a principios de la década de los 2000, ASP.NET Web Forms fue fundamental durante muchos años para la creación de sistemas empresariales y gubernamentales a gran escala. Como explica el ingeniero Yidong Zhu (2022), esta tecnología nació con el objetivo de simplificar drásticamente la programación web, permitiendo a los desarrolladores arrastrar y soltar elementos visuales en la pantalla para mostrar bases de datos o crear formularios complejos en pocos minutos.\cite{zhu2022diseno}

Sin embargo, para lograr esta simplicidad en una época donde los diseños monolíticos eran la norma, Web Forms introdujo un mecanismo llamado View State. Este componente consistía en un campo oculto que guardaba el estado de la página y se enviaba de ida y vuelta al servidor en cada clic. Tomáš Herceg (2024) señala que, aunque esto facilitaba las cosas al programador, con el tiempo volvía a las páginas lentas, propensas a parpadeos y generaba un código muy extenso que mezclaba la lógica del negocio con la parte visual, haciendo que el mantenimiento fuera una tarea pesada.\cite{herceg2024modernizing}

Además, la documentación oficial de .NET (2026), estas aplicaciones dependían obligatoriamente del sistema operativo Windows y del servidor Internet Information Services (IIS). Al ser una arquitectura tan amarrada a Windows, Microsoft decidió no incluir soporte para Web Forms en sus nuevas plataformas, obligando a que cualquier modernización requiera una reescritura completa de las interfaces.\cite{MicrosoftLibraryChangeRules}
\subsubsection{Nueva Generación de la Plataforma .NET (.NET Core y .NET moderno)}
\label{sec:nueva_generacion_net}

\subsection{Patrones de Diseño y Arquitectura de Software}
\label{sec:patrones_diseno}

\subsubsection{Patrones Arquitectónicos}
\label{sec:patrones_arquitectonicos}

\paragraph{Patrón Modelo-Vista-Controlador (MVC)} ~\\
\label{sec:mvc}
\subsection{Seguridad, Autenticación e Integración de Datos}
\label{sec:seguridad_autenticacion}

\subsubsection{Control de Acceso Basado en Roles (RBAC)}
\label{sec:rbac}

\subsubsection{Integración de Sistemas mediante APIs}
\label{sec:integracion_sistemas}

\subsubsection{Persistencia y Mapeo de Datos}
\label{sec:persistencia_mapeo}

\subsection{Propuesta de Solución}
\label{sec:propuesta_solucion}
